% ===== 底本 第 3 章《环簇 Fano 集的 SYZ 镜像对称》；片段 tr04a；源块 B132--B139 =====

% [源 B132]
\subsection{环簇 Fano 集的 SYZ 镜像构造}

% [源 B133]
在另一侧的同调镜像对称得到确认之后，我们就可以着手环簇 Fano 集的 SYZ 镜像构造了。正如上一节所解释的，紧化 $\left(\mathbb{C}^{\ast}\right)^n$ 会导致所有拉格朗日环面上的 Floer $A_{\infty}$-代数发生形变，而本节的任务就是理解这一 $A_{\infty}$-形变所带来的效果。

% [Tutor Note: B133 所说的「形变」指：在 $(\mathbb{C}^{\ast})^n\cong T^{\ast}\mathbb{T}^n$ 上 $\pi_2=0$，所有以 $L_b$ 为边界的伪全纯圆盘都是平凡的，故 Floer 代数与 de Rham 代数同构；紧化之后出现了与环簇边界除子 $D$ 相交的圆盘，$\mathfrak{m}^d$ 里便多出按圆盘类 $\beta$ 求和的那一项，这正是「形变」。可与本章上一节「$CF^{\ast}(L_b,L_b)$ 是 $CF^{\ast}(L_0,L_0)$ 的形变」一段对读。]

% [源 B134]
第一个结果是：对于环簇紧化，伪全纯圆盘的 Maslov 指标是可以完全弄清楚的：

% [源 B135]
\begin{proposition*}[Cho--Oh '06]
设 $X_{\Sigma}$ 是环簇 Fano 流形，其环簇边界除子为 $D=\displaystyle\bigcup_i D_i$，其中各 $D_i$ 是不可约分支。则对任何伪全纯圆盘 $u\colon (D^2,\partial D^2)\to (X_{\Sigma},L_b)$（其边界落在拉格朗日环面 $L_b\subseteq\left(\mathbb{C}^{\ast}\right)^n$ 上），其 Maslov 指标为
\[
  \mu (u)=2\sum_i (u\cdot D_i),
\]
其中 $u\cdot D_i$ 是 $u$ 与除子 $D_i$ 的相交数。
\end{proposition*}

% [Tutor Note: B135 公式 $\mu(u)=2\sum_i(u\cdot D_i)$ 中的系数 $2$ 来自 Maslov 指标的标准化约定：绕边界 $\partial D^2$ 走一周贡献 Maslov 指标 $2$，于是每次与环簇边界除子横截相交贡献 $2$。可核对的最简情形是 $n=1$、$X_{\Sigma}=\mathbb{P}^1$：取 $L_b=\{\lvert z\rvert=r\}$ 与圆盘 $\{\lvert z\rvert\leq r\}$，它与 $D_1=\{0\}$ 恰相交一次、与 $D_2=\{\infty\}$ 不交，故 $\mu(u)=2$，正是 B142 中 $D(v_1)$ 的情形。]

% [源 B136]
\begin{proof}
注意 Maslov 指标的计算是纯拓扑问题：在光滑同痕意义下，只有 $u$ 的像与 $D$ 的相交点会对 Maslov 指标有贡献，因此我们可以考虑最简单的模型：设有一个圆盘 $\mathbb{D}^2\subseteq\mathcal{N}_X\left(\bigcap_{i\in I} D_i\right)$，它在原点处与 $D_I\colon =\displaystyle\bigcap_{i\in I} D_i$ 相交、重数为 $\ell$，并且边界落在 $D_I$ 上的环面丛中。对这个局部模型，我们知道 $D_I$ 可以写成 $\displaystyle\left\{\prod_{i\in I}z_i=0\right\}$，于是可以直接计算（例如参见 [MS04] 的附录 D），并验证这样的圆盘的 Maslov 指标为 $2\ell\vert I\vert$。
\end{proof}

% [Tutor Note: B136 的局部模型给出 $2\ell\vert I\vert$，与 B135 的 $2\sum_i(u\cdot D_i)$ 相互印证：取 $I$ 为单点集且 $\ell=1$，即得 $\mu=2$；若 $I$ 含 $m$ 个分支、每个重数 $1$，则得 $\mu=2m$。又：源文写 $\mathbb{D}^2\subseteq\mathcal{N}_X(\bigcap_{i\in I}D_i)$，$\mathcal{N}_X(\cdot)$ 是 $\bigcap_{i\in I}D_i$ 在 $X$ 中的法丛，此处「$\subseteq$」应理解为「圆盘落在该法丛的管状邻域中」，是原文的一处记号混用，译文照原样保留。]

% [源 B137]
特别地，存在 Maslov 指标为 $2$ 的伪全纯圆盘，它们对 $A_{\infty}$-结构的曲率项 $\mathfrak{m}^0$ 有贡献，从而可能妨碍拉格朗日环面的 Floer 上同调成为良定义的。为了计算这一项，我们需要对这些伪全纯圆盘进行分类。这正是 [CO06] 的下述主要定理：

% [Tutor Note: B137 把 $\mathfrak{m}^0$ 称为「曲率项」（curvature term）：它是 $A_{\infty}$ 结构中唯一没有输入、只产生一个输出的运算，一旦非零，$\mathfrak{m}^1$ 一般就不再有 $\left(\mathfrak{m}^1\right)^2=0$。B148--B154 将证明：只要 $L_b$ 不围出 Maslov 指标 $\leq 0$ 的全纯圆盘，$(\mathfrak{m}^1)^2=0$ 仍然成立——这解释了本节为何先显式计算 $\mathfrak{m}^0$、再回来谈 Floer 上同调的良定义。]

% [源 B138]
\begin{theorem*}[Cho--Oh '06]
设 $\tilde{L}\subseteq\mathbb{C}^N\setminus Z(\Sigma)$ 是实 $N$-环面 $\mathbb{T}^N$ 的一个固定轨道。任何全纯映射 $u\colon (\mathbb{D}^2,\partial\mathbb{D}^2)\to (X_{\Sigma},L)$ 都可以提升为全纯映射
\[
  \tilde{u}\colon (\mathbb{D}^2,\partial\mathbb{D}^2)\to (\mathbb{C}^N\setminus Z(\Sigma),\tilde{L}),
\]
使得每个齐次坐标函数 $z_1 (\tilde{u}),\dotsb, z_N (\tilde{u})$ 都由带常数因子的 Blaschke 乘积给出，即
\[
  z_i (\tilde{u})=c_i\cdot\prod_{j=1}^{\mu_i}\frac{z-\alpha_{i,j}}{1-\bar{\alpha}_{i,j} z},
\]
其中 $c_i\in\mathbb{C}^{\ast}$，$\mu_i\in\mathbb{Z}_{>0}$，$\alpha_{i,j}\in\operatorname{Int} (D^2)$。任何两个这样的提升都由 $\tilde{u}'=t\cdot\tilde{u}$ 相联系，其中 $t\in T(\Sigma)$ 是某个常数元素。
\end{theorem*}

% [Tutor Note: B138 的「提升」把 $X_{\Sigma}$ 上的全纯圆盘拉到 $\mathbb{C}^N\setminus Z(\Sigma)$ 上；每个分量是 Blaschke 乘积这件事，本质上是「边界落在 $\tilde{L}$ 上」迫使 $\lvert z_i(\tilde{u})\rvert$ 在边界上为常数：当 $\lvert z\rvert=1$ 时 $\left\lvert\frac{z-\alpha}{1-\bar{\alpha}z}\right\rvert=1$，故每个因子只改变圆盘内部的零点位置、不改变边界模长。于是 $z_i(\tilde{u})$ 在圆盘内的零点个数 $\mu_i$ 是拓扑不变量，B139 据此把它等同于与 $D_i$ 的相交重数。]

% [Tutor Note: B138 中的 $T(\Sigma)$ 只被原文称为「某个常数元素」（constant element），未给出进一步定义，此处照译；它的作用是刻画提升的不唯一性（两个提升相差一个常值环面元素）。同时提醒：$N$ 是环簇边界除子的个数（即扇的射线数），不要与 $n=\dim X_{\Sigma}$ 混淆——B145 的求和指标跑到 $N$，而 $b\in\mathbb{R}^n$。源文 $D^2$（B135、B138）与 $\mathbb{D}^2$（B136--B138）混用，均指单位圆盘，转写时照录。]

% [Tutor Note: B138 的 $\mu_i\in\mathbb{Z}_{>0}$ 疑当作 $\mathbb{Z}_{\geq 0}$：B142 的分类中 $\tilde{u}(z)=(c_1,\dotsc,c_{i-1},c_i z,c_{i+1},\dotsc,c_N)$ 只有一个分量含零点（$\mu_i=1$，其余 $\mu_j=0$），若要求所有 $\mu_i>0$，则 B142 的圆盘不满足 B138 的结论。原文如此，转写照录，未改。]

% [源 B139]
特别地，这些数 $\mu_i$ 就是伪全纯圆盘 $u(\mathbb{D}^2)$ 与环簇边界分支 $D_i$ 的相交重数。利用这一分类，我们就可以完全计算具有任意 Maslov 指标的伪全纯圆盘对 $\mathfrak{m}^0$ 的贡献。

% [Tutor Note: 由 B135 与 B138--B139 得 $\mu(u)=2\sum_i(u\cdot D_i)=2\sum_i\mu_i$，即 Maslov 指标只由 Blaschke 因子的总个数决定；又 $\mu_i\in\mathbb{Z}_{\geq 0}$（见 B138 补注），故 $\mu(u)=2$ 当且仅当恰有一个 $i$ 使 $\mu_i=1$、其余 $\mu_j=0$——这正是 B142 中「只与一个边界分支相交一次」的圆盘，可作为核对 B142 的分类是否完备的锚点。]

% [Tutor Note: B139 说可以「完全计算任意 Maslov 指标圆盘的贡献」，而 B140 的引理随即指出 Maslov 指标 $2k$（$k\geq 2$）的求和整体为 $0$。也就是说，要证明消失也必须先有 B138 的分类；只分类 Maslov 指标 $2$ 的圆盘是不够的。这一逻辑顺序在翻译时保留原文的「先分类、再消失」两段式。]
